% =====================================================================
%  5.3 Sinussetningen — fasit
% =====================================================================

\fk{Del 1 — Sinussetningen og beviset}
\fo{1.1} a) $60^\circ$ \ b) $50^\circ$ \ c) $53^\circ$ \ d) $50^\circ$
\fo{1.2} a) $\frac{b}{\sin 30^\circ} = \frac{6\sqrt3}{\sin 60^\circ}$ \ b) $b = 6\sqrt3 \cdot \frac12 \cdot \frac{2}{\sqrt3} = 6$
\fo{1.3} $C = 70^\circ$, $b \approx 7,9$, $c \approx 8,6$
\fo{1.4} a) $b = \frac{12,0\sin 80^\circ}{\sin 35^\circ} \approx 20,6$ \ b) $A = 61^\circ$, $a = \frac{9,5\sin 61^\circ}{\sin 71^\circ} \approx 8,8$ \ c) $a = \frac{5,0\sin 112^\circ}{\sin 28^\circ} \approx 9,9$
\fo{1.5} $a = \frac{6\sqrt2 \cdot \sin 120^\circ}{\sin 45^\circ} = 6\sqrt2 \cdot \frac{\sqrt3}{2} \cdot \frac{2}{\sqrt2} = 6\sqrt3$. Videre er $C = 180^\circ - 120^\circ - 45^\circ = 15^\circ$.
\fo{1.6} a) $180^\circ - 62^\circ - 71^\circ = 47^\circ$ \ b) $PD = \frac{100\sin 71^\circ}{\sin 47^\circ} \approx 129,3$ m (overfor $Q$) og $QD = \frac{100\sin 62^\circ}{\sin 47^\circ} \approx 120,7$ m \ c) $129,3 \cdot \sin 62^\circ \approx 114$ m
\fo{1.7} Feilen er overgangen til $c = \frac{9,0\cdot\sin 40^\circ}{\sin 75^\circ}$: brøken er snudd. Riktig er $c = \frac{9,0 \cdot \sin 75^\circ}{\sin 40^\circ} \approx 13,5$. Kontroll: $C > A$, så $c$ skal være større enn $a = 9,0$.
\fo{1.8} a) Alltid. $\sin A = \sin B$ gir $\frac{\sin A}{a} = \frac{\sin A}{b}$, så $a = b$.
b) Aldri. Da er ingen vinkel med motstående side kjent. Hver likning har da to ukjente, for eksempel $\frac{\sin A}{a} = \frac{\sin B}{b}$ med både $A$ og $B$ ukjent. Cosinussetningen (5.4) trengs.
c) Alltid. Vi finner tredje vinkel med vinkelsummen. Da er alle vinklene kjent, og siden mellom de to vinklene er et kjent par. Sinussetningen gir resten.
\fo{1.9} a) Del begge sider på $\frac12 abc$. Venstre side blir $\frac{\frac12 bc\sin A}{\frac12 abc} = \frac{\sin A}{a}$, siden $\frac12$, $b$ og $c$ forkortes. Høyre side blir $\frac{\frac12 ac\sin B}{\frac12 abc} = \frac{\sin B}{b}$, siden $\frac12$, $a$ og $c$ forkortes. Likheten består, fordi vi deler like store uttrykk på det samme tallet.
b) $\frac{\sin 90^\circ}{a} = \frac{\sin B}{b}$ og $\sin 90^\circ = 1$ gir $\frac1a = \frac{\sin B}{b}$, altså $\sin B = \frac ba$. Her er $a$ hypotenus og $b$ motstående katet til $B$. Det er definisjonen $\sin B = \frac{\text{motstående}}{\text{hypotenus}}$.
\fo{1.10} Ett svar: $A = 30^\circ$, $B = 60^\circ$, $C = 90^\circ$. Da er $\frac ba = \frac{\sin 60^\circ}{\sin 30^\circ} = \sqrt3$. Et annet: $A = 30^\circ$, $B = 120^\circ$, $C = 30^\circ$, fordi $\sin 120^\circ = \sin 60^\circ$. Flere svar er riktige.
\fo{Sjekk 1} a) $60^\circ$ \ b) $b = \frac{3\sin 90^\circ}{\sin 30^\circ} = 6$ \ c) Vi trenger en vinkel og siden overfor for å bruke sinussetningen. Vinkelsummen gir vinkelen overfor den kjente siden.

\fk{Del 2 — Det tvetydige tilfellet}
\fo{2.1} a) $30^\circ$ og $150^\circ$ \ b) $90^\circ$ \ c) ingen, siden $\sin v \le 1$ \ d) $60^\circ$ og $120^\circ$
\fo{2.2} a) $\sin B = \frac{6\sqrt3 \cdot \sin 30^\circ}{6} = \frac{\sqrt3}{2}$ \ b) $B = 60^\circ$ eller $B = 120^\circ$ \ c) $30^\circ + 120^\circ = 150^\circ < 180^\circ$, så begge gir trekant. $B = 60^\circ$: $C = 90^\circ$ og $c = 12$. $B = 120^\circ$: $C = 30^\circ$ og $c = 6$.
\fo{2.3} a) $\sin B \approx 0,750$ \ b) $B_1 \approx 48,6^\circ$ og $B_2 \approx 131,4^\circ$ \ c) $40^\circ + 131,4^\circ = 171,4^\circ < 180^\circ$. Ja, begge gir trekant.
\fo{2.4} $\sin B = \frac{9,0\sin 35^\circ}{7,0} \approx 0,737$. To trekanter, siden $35^\circ + 132,5^\circ < 180^\circ$.
Trekant 1: $B \approx 47,5^\circ$, $C \approx 97,5^\circ$, $c \approx 12,1$.
Trekant 2: $B \approx 132,5^\circ$, $C \approx 12,5^\circ$, $c \approx 2,6$.
\fo{2.5} a) Én trekant. $\sin B = \frac{6,0\sin 62^\circ}{9,0} \approx 0,589$, $B \approx 36,1^\circ$. Den andre kandidaten $143,9^\circ$ gir sum over $180^\circ$ med $62^\circ$. $C \approx 81,9^\circ$ og $c \approx 10,1$.
b) Ingen trekant. $\sin B = \frac{7,0 \sin 40^\circ}{3,0} \approx 1,50 > 1$.
c) Én trekant. $\sin B \approx 0,658$, $B \approx 41,1^\circ$. Den andre kandidaten $138,9^\circ$ er for stor sammen med $110^\circ$. $C \approx 28,9^\circ$ og $c \approx 5,1$.
\fo{2.6} a) $A$ og $C$ er faste. Tegn en sirkel med radius $6,0$ km om $C$. Den skjærer kystlinja i to punkter. b) $\sin B = \frac{8,0 \sin 25^\circ}{6,0} \approx 0,563$, så $B \approx 34,3^\circ$ eller $B \approx 145,7^\circ$. Begge går ($25^\circ + 145,7^\circ < 180^\circ$). Avstand fra havna: $\approx 12,2$ km eller $\approx 2,3$ km.
\fo{2.7} Feilen er overgangen til «$B \approx 63,3^\circ$ (eneste løsning)». Også $B = 180^\circ - 63,3^\circ = 116,7^\circ$ har $\sin B \approx 0,893$, og $44^\circ + 116,7^\circ < 180^\circ$. Riktig svar: to trekanter.
Trekant 1: $B \approx 63,3^\circ$, $C \approx 72,7^\circ$, $c \approx 9,6$.
Trekant 2: $B \approx 116,7^\circ$, $C \approx 19,3^\circ$, $c \approx 3,3$.
\fo{2.8} $\sin B = \frac{9,0 \sin A}{6,0} = 1,5\sin A$.
Ingen trekant når $1,5 \sin A > 1$, altså $\sin A > \frac23$, dvs. $41,8^\circ < A < 138,2^\circ$: da er $\sin B > 1$. Når $A \ge 138,2^\circ$ er $1,5\sin A > \sin A = \sin(180^\circ - A)$, så den spisse $B$ er større enn $180^\circ - A$. Da er $A + B > 180^\circ$, og det finnes ingen trekant.
Én trekant: bare når $\sin B = 1$, dvs. $A \approx 41,8^\circ$ og $B = 90^\circ$.
To trekanter når $A < 41,8^\circ$. Da er $\sin B < 1$, og $1,5\sin A > \sin A$ gir $B > A$. Derfor er $A + (180^\circ - B) < 180^\circ$, og $A + B < 180^\circ$ siden begge er spisse. Med svingarmen: $h = 9,0\sin A < 6,0 < 9,0$.
\fo{2.9} a) Alltid. Er $A$ stump, må $B$ være spiss for at summen skal bli under $180^\circ$. Da gir $180^\circ - B$ en stump $B$ og sum over $180^\circ$. Det finnes altså høyst én trekant, og noen ganger ingen.
b) Noen ganger. $A = 30^\circ$, $b = 10$: $a = 7$ gir to trekanter, men $a = 3$ gir ingen, fordi $a < h = 5$.
\fo{2.10} Eksempel med $A = 30^\circ$ og $b = 10$. To trekanter: $a = 5\sqrt2$ gir $\sin B = \frac{10 \cdot \frac12}{5\sqrt2} = \frac{1}{\sqrt2} = \frac{\sqrt2}{2}$, så $B = 45^\circ$ eller $B = 135^\circ$. Begge har sum under $180^\circ$ ($30^\circ + 135^\circ = 165^\circ$). Én trekant: $a = 5$ gir $\sin B = 1$, så $B = 90^\circ$ og ingen annen vinkel. Flere svar er riktige.
\fo{Sjekk 2} a) $36,9^\circ$ og $143,1^\circ$ \ b) Ingen, siden $h = 12\sin 30^\circ = 6 > 5$ \ c) Fordi to ulike vinkler har samme sinus. Er summen med den stumpe vinkelen $180^\circ$ eller mer, finnes ikke trekanten.

\fk{Ekstra utfordring}
\fo{U1} $b = \frac{4\sqrt2 \cdot \sin 60^\circ}{\sin 45^\circ} = \frac{4\sqrt2 \cdot \frac{\sqrt3}{2}}{\frac{\sqrt2}{2}} = 4\sqrt3$. Videre er $C = 75^\circ$.
\fo{U2} a) $C = 61^\circ$ og $AC = \frac{70\sin 71^\circ}{\sin 61^\circ} \approx 75,7$ m \ b) Bredden er høyden fra $C$: $75,7 \cdot \sin 48^\circ \approx 56$ m.
\fo{U3} $C = 65^\circ$, $b = \frac{12 \sin 65^\circ}{\sin 50^\circ} \approx 14,2$ og $T = \frac12 \cdot 12 \cdot 14,2 \cdot \sin 65^\circ \approx 77,2$.
\fo{U4} Alltid. Anta $A > B$. Er begge spisse, vokser sinus på $[0^\circ, 90^\circ]$, så $\sin A > \sin B$. Er $A = 90^\circ$, er $\sin A = 1$ og $\sin B < 1$. Er $A$ stump, er $B < 180^\circ - A$ som er spiss, og $\sin B < \sin(180^\circ - A) = \sin A$. Sinussetningen gir $\frac ab = \frac{\sin A}{\sin B} > 1$, altså $a > b$.
\fo{U5} a) Utskriften er
\lstinline[style=laereverk,numbers=none,frame=none]|2 (0, [])|,
\lstinline[style=laereverk,numbers=none,frame=none]|3 (1, [90.0])|,
\lstinline[style=laereverk,numbers=none,frame=none]|4 (2, [48.6, 131.4])|,
\lstinline[style=laereverk,numbers=none,frame=none]|8 (1, [22.0])| og
\lstinline[style=laereverk,numbers=none,frame=none]|(1, [32.5])|.
$a = 2$: ingen ($h = 3$). $a = 3$: én, $B = 90^\circ$. $a = 4$: to. $a = 8$: én, siden $a \ge b$. Siste linje er $\mathrm{ssv}(7, 4, 70)$: også $a > b$, så én.
b) Ta $a = 4$. Da er $B_1 = 48,6^\circ$ og $180^\circ - B_1 = 131,4^\circ$ er også en løsning av $\sin B = 0,75$. Uten testen ville programmet bare gitt én trekant, selv om det finnes to.
c) Grunnen til nøyaktig én trekant er $a > b$: da er $B < A$, og $180^\circ - B_1$ gir sum over $180^\circ$. Eksempel: \lstinline[style=laereverk,numbers=none,frame=none]|ssv(10, 7, 110)| gir \lstinline[style=laereverk,numbers=none,frame=none]|(1, [41.1])|. Med $a < b$ og stump $A$ finnes ingen: \lstinline[style=laereverk,numbers=none,frame=none]|ssv(6, 9, 150)| gir \lstinline[style=laereverk,numbers=none,frame=none]|(0, [])|, siden $150^\circ + 48,6^\circ > 180^\circ$ og $150^\circ + 131,4^\circ > 180^\circ$.
\fo{U6} a) Vinkel $C$ spenner over buen $BD$. Siden $BD$ er en diameter, er buen en halv sirkel, $180^\circ$. Periferivinkelen er halvparten: $90^\circ$.
b) $BD = 2R$ er hypotenusen i den rettvinklede trekanten $BCD$. Siden motstående til $D$ er $BC = a$.
c) Hypotenusen er $2R = 13$, så $R = \frac{13}{2}$. Sirkelens sentrum ligger midt på hypotenusen.
d) Er $A$ stump, ligger $A$ og $D$ på hver sin side av $BC$. De spenner over hver sin bue, og buene er til sammen $360^\circ$. Da er $\angle A + \angle D = 180^\circ$, så $\angle D = 180^\circ - \angle A$ og $\sin D = \sin A$. Resten av beviset er likt.
